\label{anlage:user-stories}

Aus der Recherche und den Interviews konsolidierte User Stories der vier zentralen Akteursgruppen, mit der daraus abgeleiteten Systemarchitektur und dem Anforderungsmapping, das die User Stories in funktionale Anforderungscluster übersetzt.

\begin{Table}[title=Übersicht]
\begin{SemioTable}{@{}T{0.7\linewidth-2\tabcolsep}U{0.3\linewidth-2\tabcolsep}@{}}
\SemioTableHeaderRow{Akteursgruppe & Konsolidierte User Stories}
\SemioTableRow{Mitarbeiter:in einer Bauteilbörse & 10}
\SemioTableRow{Architekt:in & 12}
\SemioTableRow{Tragwerksplaner:in & 6}
\SemioTableRow{Energieberater:in & 7}
\SemioTableRow{\textbf{Gesamt} & \textbf{35}}
\end{SemioTable}
\end{Table}

\appendixsection[S]{User Stories}

\subsection{Mitarbeiter:in einer Bauteilbörse}

\paragraph{US-BB-01 — Baukomponentendaten erfassen} Als Mitarbeiter:in einer Bauteilbörse möchte ich wiederverwendbare Baukomponenten mit relevanten technischen, logistischen und qualitativen Informationen erfassen, damit Planungsteams sie als belastbare Grundlage für Entwurfs- und Entscheidungsprozesse nutzen können.

\paragraph{US-BB-02 — Baukomponenten zu planungsrelevanten Einheiten bündeln} Als Mitarbeiter:in einer Bauteilbörse möchte ich einzelne Baukomponenten zu sinnvollen Gruppen oder Paketen zusammenfassen, damit Planungsteams mit verwendbaren Mengen und zusammenhängenden Materialbeständen arbeiten können.

\paragraph{US-BB-03 — Verfügbarkeit transparent machen} Als Mitarbeiter:in einer Bauteilbörse möchte ich zeitliche Verfügbarkeiten von Baukomponenten transparent darstellen, damit Rückbau, Lagerung, Planung und Wiedereinbau besser koordiniert werden können.

\paragraph{US-BB-04 — Status von Baukomponenten verwalten} Als Mitarbeiter:in einer Bauteilbörse möchte ich den aktuellen Status von Baukomponenten verwalten, damit Planungsteams erkennen können, welche Elemente verfügbar, reserviert, in Prüfung oder nicht mehr nutzbar sind.

\paragraph{US-BB-05 — Bestandsdaten standardisieren} Als Mitarbeiter:in einer Bauteilbörse möchte ich Informationen aus Gebäudebeständen in einem einheitlichen Datenformat erfassen, damit sie vergleichbar, auswertbar und in weiteren Planungsschritten nutzbar sind.

\paragraph{US-BB-06 — Kontextinformationen dokumentieren} Als Mitarbeiter:in einer Bauteilbörse möchte ich ergänzende Informationen zum ursprünglichen Einsatz, zur gestalterischen Absicht und zum Nutzungskontext einer Baukomponente dokumentieren, damit Planungsteams deren Wiederverwendung qualifiziert bewerten können.

\paragraph{US-BB-07 — Baukomponentenhistorie nachvollziehbar machen} Als Mitarbeiter:in einer Bauteilbörse möchte ich die Herkunft und Nutzungsgeschichte einer Baukomponente dokumentieren, damit Planungsteams deren Qualität, Eignung und potenzielle Risiken besser einschätzen können.

\paragraph{US-BB-08 — Informationslücken sichtbar machen} Als Mitarbeiter:in einer Bauteilbörse möchte ich fehlende oder noch zu prüfende Informationen sichtbar machen, damit offene Nachweise gezielt an zuständige Fachplaner:innen übergeben werden können.

\paragraph{US-BB-09 — Kostenstruktur transparent darstellen} Als Mitarbeiter:in einer Bauteilbörse möchte ich die relevanten Kostenbestandteile wiederverwendbarer Baukomponenten transparent darstellen, damit Planungsteams wirtschaftliche Auswirkungen realistisch bewerten können.

\paragraph{US-BB-10 — Zusatzleistungen modular anbieten} Als Mitarbeiter:in einer Bauteilbörse möchte ich projektbezogene Zusatzleistungen modular anbieten, damit Planungsteams je nach Projektphase und Informationsbedarf passende Unterstützungsleistungen auswählen können.

\subsection{Architekt:in}

\paragraph{US-A-01 — Verfügbare Baukomponenten frühzeitig einsehen} Als Architekt:in möchte ich verfügbare wiederverwendbare Baukomponenten bereits in frühen Planungsphasen einsehen, damit der Entwurf von Beginn an auf realen Materialverfügbarkeiten aufbauen kann.

\paragraph{US-A-02 — Mit Baukomponenten entwerfen} Als Architekt:in möchte ich Baukomponenten anhand relevanter Eigenschaften suchen und filtern, damit ich Entwurfskonzepte mit vorhandenen Elementen entwickeln und prüfen kann.

\paragraph{US-A-03 — Baukomponenten für Entwurfsvarianten sichern} Als Architekt:in möchte ich relevante Baukomponenten temporär für bestimmte Entwurfsvarianten sichern können, damit die Planung auf einer verlässlicheren Materialverfügbarkeit basiert.

\paragraph{US-A-04 — Entwurfsvarianten vergleichen} Als Architekt:in möchte ich unterschiedliche Entwurfsvarianten vergleichen, damit gestalterische, technische, ökologische, wirtschaftliche und organisatorische Auswirkungen fundiert bewertet werden können.

\paragraph{US-A-05 — Baukomponenten digital in den Entwurf integrieren} Als Architekt:in möchte ich wiederverwendbare Baukomponenten als strukturierte digitale Planungsobjekte übernehmen, damit sie direkt in Entwurfs- und Modellierungsprozesse eingebunden werden können.

\paragraph{US-A-06 — Entwurfsrelevante Risiken erkennen} Als Architekt:in möchte ich entwurfsrelevante Risiken und fehlende Informationen zu Baukomponenten erkennen, damit ich deren Eignung als Planungsgrundlage realistisch einschätzen kann.

\paragraph{US-A-07 — Änderungen an geplanten Baukomponenten nachvollziehen} Als Architekt:in möchte ich über relevante Änderungen an eingeplanten Baukomponenten informiert werden, damit ich Entwurf, Termine und Abstimmungen rechtzeitig anpassen kann.

\paragraph{US-A-08 — Gestalterische Eigenschaften berücksichtigen} Als Architekt:in möchte ich ästhetische und materielle Eigenschaften wiederverwendbarer Baukomponenten gezielt berücksichtigen, damit Wiederverwendung als bewusster Bestandteil des Entwurfs eingesetzt werden kann.

\paragraph{US-A-09 — Aufgaben zur Wiederverwendung phasenbezogen steuern} Als Architekt:in möchte ich wiederverwendungsbezogene Aufgaben entlang der Planungsphasen strukturieren, damit Entscheidungen, Prüfungen und Dokumentationen zum richtigen Zeitpunkt erfolgen.

\paragraph{US-A-10 — Risiken systematisch bewerten} Als Architekt:in möchte ich Risiken zu geplanten Baukomponenten systematisch angezeigt bekommen, damit ich ihre Auswirkungen auf Entwurf, Kosten, Termine und Umsetzbarkeit bewerten kann.

\paragraph{US-A-11 — Baukomponenten flexibel filtern} Als Architekt:in möchte ich Baukomponenten über frei kombinierbare Kriterien filtern, damit ich geeignete Elemente gezielt für unterschiedliche Entwurfsanforderungen auswählen kann.

\paragraph{US-A-12 — Wiederverwendungsprozesse den Planungsphasen zuordnen} Als Architekt:in möchte ich wiederverwendungsbezogene Entscheidungen und Aufgaben den jeweiligen Planungsphasen zuordnen, damit der Wiederverwendungsprozess nachvollziehbar in den Planungsablauf integriert wird.

\subsection{Tragwerksplaner:in}

\paragraph{US-T-01 — Tragwerksrelevante Baukomponenten identifizieren} Als Tragwerksplaner:in möchte ich wiederverwendbare Baukomponenten anhand strukturell relevanter Eigenschaften suchen und filtern, damit ich geeignete Elemente für tragende Anwendungen schnell identifizieren kann.

\paragraph{US-T-02 — Fehlende Tragwerksinformationen erkennen} Als Tragwerksplaner:in möchte ich fehlende oder unvollständige Tragwerksinformationen erkennen, damit notwendige Prüfungen, Annahmen und Nachweise definiert werden können.

\paragraph{US-T-03 — Frühes fachliches Feedback geben} Als Tragwerksplaner:in möchte ich frühzeitig Rückmeldungen zu vorgeschlagenen Anordnungen von Baukomponenten geben, damit tragwerksrelevante Fragen vor zentralen Entwurfsentscheidungen geklärt werden können.

\paragraph{US-T-04 — Konstruktive Lösungsprinzipien dokumentieren} Als Tragwerksplaner:in möchte ich mögliche konstruktive Lösungs- und Anschlussprinzipien dokumentieren, damit wiederverwendbare Baukomponenten realistisch in die weitere Planung integriert werden können.

\paragraph{US-T-05 — Tragwerksvarianten vergleichen} Als Tragwerksplaner:in möchte ich unterschiedliche Tragwerksvarianten vergleichen, damit technische, wirtschaftliche, ökologische und gestalterische Auswirkungen fundiert gegeneinander abgewogen werden können.

\paragraph{US-T-06 — Nachweisfähige Informationen exportieren} Als Tragwerksplaner:in möchte ich relevante Berechnungen, Annahmen, Prüfresultate und Entscheidungen standardisiert exportieren, damit wiederverwendbare Baukomponenten in Genehmigungs-, Ausschreibungs- und Ausführungsprozesse überführt werden können.

\subsection{Energieberater:in}

\paragraph{US-E-01 — Frühe Umweltwirkung abschätzen} Als Energieberater:in möchte ich relevante Material-, Herkunfts- und Logistikinformationen zu wiederverwendbaren Baukomponenten erhalten, damit ökologische Auswirkungen bereits in frühen Planungsphasen abgeschätzt werden können.

\paragraph{US-E-02 — Wiederverwendbare Baukomponenten in Bewertungen integrieren} Als Energieberater:in möchte ich wiederverwendbare Baukomponenten mit bewertungsrelevanten Kennwerten verknüpfen, damit unterschiedliche Material- und Konstruktionsvarianten vergleichbar analysiert werden können.

\paragraph{US-E-03 — Ökologische Wirkung von Varianten vergleichen} Als Energieberater:in möchte ich die ökologische Wirkung verschiedener Entwurfsvarianten vergleichen, damit Planungsteams die Auswirkungen von Materialentscheidungen nachvollziehen können.

\paragraph{US-E-04 — Varianten energetisch prüfen} Als Energieberater:in möchte ich Varianten mit relevanten Modellinformationen verknüpfen, damit deren energetische Auswirkungen ohne unverhältnismäßigen Modellierungsaufwand geprüft werden können.

\paragraph{US-E-05 — Grenzwerte und Zielwerte prüfen} Als Energieberater:in möchte ich Projekte mit relevanten ökologischen und energetischen Ziel- oder Grenzwerten vergleichen, damit Wiederverwendungsentscheidungen messbar bewertet werden können.

\paragraph{US-E-06 — Optimierungsmöglichkeiten aufzeigen} Als Energieberater:in möchte ich erkennen, wann wiederverwendbare Baukomponenten durch ergänzende Maßnahmen verbessert werden sollten, damit zirkuläre und energetische Anforderungen gemeinsam erfüllt werden können.

\paragraph{US-E-07 — Bewertungsergebnisse exportieren} Als Energieberater:in möchte ich relevante Kennwerte, Diagramme und Bewertungsergebnisse exportieren, damit Auftraggeber:innen, Planungsteams und Behörden die Entscheidungsgrundlagen nachvollziehen können.

\appendixsection[SA]{Systemarchitektur für das Anforderungsmapping}

Die Systemstruktur wird aus den konsolidierten User Stories abgeleitet. Die wiederkehrenden Anforderungen verweisen auf zwei Hauptbereiche: die Einspeiseplattform mit dem Bauteilportal und das Design Tool mit den Teilbereichen Katalog, Entwurf und Konfigurator. Die folgende Tabelle dient als Leselogik für das anschließende Anforderungsmapping.

\begin{Table}[title=Systemarchitektur]
\begin{SemioTable}{@{}T{0.16\linewidth-2\tabcolsep}T{0.16\linewidth-2\tabcolsep}T{0.20\linewidth-2\tabcolsep}T{0.48\linewidth-2\tabcolsep}@{}}
\SemioTableHeaderRow{Hauptbereich & Teilbereich & Hauptfunktion & Einordnung}
\SemioTableRow{Einspeiseplattform & Bauteilportal & Datenbereitstellung & Bündelt Anforderungen zur Erfassung, Qualifizierung und Bereitstellung von Baukomponentendaten.}
\SemioTableRow{Design Tool & Tool (Katalog) & Komponentenauswahl & Bündelt Anforderungen zur Sichtbarkeit, Suche, Filterung, Vergleichbarkeit und fachlichen Einschätzung von Baukomponenten.}
\SemioTableRow{Design Tool & Tool (Entwurf) & Entwurfsbearbeitung & Bündelt Anforderungen zur Verwendung, Kombination, Bewertung und Dokumentation von Baukomponenten im Entwurf.}
\SemioTableRow{Design Tool & Tool (Konfigurator) & Projektbezogene Abstimmung & Bündelt Anforderungen zur Abstimmung mit Bauteilbörsen, insbesondere zu Verfügbarkeit, Reservierung, Bestellung und Zusatzleistungen.}
\end{SemioTable}
\end{Table}

\appendixsection[AM]{Anforderungsmapping der User Stories}
\label{anlage:anforderungsmapping}

Das Mapping übersetzt die User Stories in funktionale Anforderungscluster. Die Zuordnung folgt der aus den User Stories abgeleiteten Systemstruktur.

\SemioTableLong{Anforderungsmapping}{0.10,0.40,0.26,0.24}{ID & Funktionaler Anforderungscluster & Zuordnung & Zugeordnete User Stories}{%
  \SemioTableRow{\textbf{REQ-01} & Baukomponentendaten erfassen und standardisieren & Einspeiseplattform (Bauteilportal) & US-BB-01, US-BB-05}
  \SemioTableRow{\textbf{REQ-02} & Baukomponentengruppen darstellen und auswählbar machen & Design Tool (Katalog) & US-BB-02}
  \SemioTableRow{\textbf{REQ-03} & Status-, Verfügbarkeits- und Nachweisdaten pflegen & Einspeiseplattform (Bauteilportal) & US-BB-03, US-BB-04, US-BB-08}
  \SemioTableRow{\textbf{REQ-04} & Kontext, Herkunft und Baukomponentenhistorie dokumentieren & Einspeiseplattform (Bauteilportal) & US-BB-06, US-BB-07}
  \SemioTableRow{\textbf{REQ-05} & Kosteninformationen bereitstellen & Einspeiseplattform (Bauteilportal) & US-BB-09}
  \SemioTableRow{\textbf{REQ-06} & Baukomponenten suchen, filtern und auswählen & Design Tool (Katalog) & US-A-01, US-A-02, US-A-11, US-T-01}
  \SemioTableRow{\textbf{REQ-07} & Risiken, fehlende Fachinformationen und Eignung anzeigen & Design Tool (Katalog) & US-A-06, US-T-02}
  \SemioTableRow{\textbf{REQ-08} & Gestalterische und materielle Eigenschaften bewerten & Design Tool (Katalog) & US-A-08}
  \SemioTableRow{\textbf{REQ-09} & Ökologische, energetische und leistungsbezogene Kennwerte bereitstellen & Design Tool (Katalog) & US-E-01, US-E-02}
  \SemioTableRow{\textbf{REQ-10} & Digitale Baukomponenten in den Entwurf überführen & Design Tool (Katalog) & US-A-05}
  \SemioTableRow{\textbf{REQ-11} & Baukomponenten reservieren, anfragen und servicebezogen konfigurieren & Design Tool (Konfigurator) & US-A-03, US-BB-10}
  \SemioTableRow{\textbf{REQ-12} & Entwurfs- und Tragwerksvarianten entwickeln und vergleichen & Design Tool (Entwurf) & US-A-04, US-T-05}
  \SemioTableRow{\textbf{REQ-13} & Ökologische und energetische Variantenbewertung unterstützen & Design Tool (Entwurf) & US-E-03, US-E-04, US-E-05, US-E-06}
  \SemioTableRow{\textbf{REQ-14} & Konstruktive Integration und fachliches Feedback unterstützen & Design Tool (Entwurf) & US-T-03, US-T-04}
  \SemioTableRow{\textbf{REQ-15} & Änderungen und Risiken geplanter Baukomponenten nachvollziehen & Design Tool (Entwurf) & US-A-07, US-A-10}
  \SemioTableRow{\textbf{REQ-16} & Aufgaben und Entscheidungen den Planungsphasen zuordnen & Design Tool (Entwurf) & US-A-09, US-A-12}
  \SemioTableRow{\textbf{REQ-17} & Dokumentation, Berichte und exportierbare Nachweise erzeugen & Design Tool (Entwurf) & US-T-06, US-E-07}
}
